\section{Plan de negocio}
\label{anexo:plan-negocio}

\subsection{Introducción}

\subsubsection{Considerandos}

La investigación llevada a cabo en la tesis abre la posibilidad de creación de un emprendimiento comercial ya que los hallazgos que arroja coinciden con necesidades insatisfechas detectadas en el mercado de los asistentes virtuales.

En el diseño de la propuesta de valor es crucial recordar que los datos son el activo más importante de las organizaciones. Es poco probable que una empresa permita la copia de toda la información de sus clientes en una base de datos externa y fuera de su control.

\subsubsection{Actores del mercado}

Como en toda industria existen diferentes tipos de jugadores de acuerdo a la complejidad y volumen de conversaciones que un asistente virtual debe manejar.

En el estamento más básico y de bajo volumen podemos encontrar a los proveedores de lo que se podría llamar ``haga su propio bot'', que ofrecen una solución integrada funcionando en una aplicación web. Los usuarios aprenden solos a usar el producto y gestionan el mantenimiento de su bot de manera totalmente autónoma. Es la solución ideal para dueños de tiendas online que no requieran integraciones complejas con otros sistemas.

Luego existe un grupo de oferentes de desarrollo de asistentes virtuales para los casos en los que se necesite muy alta flexibilidad y ajuste a necesidades particulares, con una llegada directa a los desarrolladores. En estos casos el volumen de conversaciones es bajo o intermedio. Este grupo es conocido como las ``agencias de bots'' o ``consultoras boutique''.

En el estamento siguiente podemos encontrar a aquellos proveedores que ofrecen una plataforma de autogestión pero para resolver conversaciones de complejidad intermedia y volumen alto. Estas empresas ofrecen horas de desarrollo de equipos propios para la puesta en marcha de los proyectos y luego los equipos internos del dueño del bot continúan con la administración y mantenimiento de manera autónoma. En vez de una tarifa plana generalmente cobran una tarifa por caso, que es lo que se conoce comúnmente como una sesión. Clientes comunes de este tipo de empresas son bancos, empresas de telecomunicaciones y de prestación de servicios de salud.

Cuando se requiere un máximo nivel de robustez el mercado acude a lo que se llaman las soluciones ``world class'', que constituyen el tope de gama en desarrollo y mantenimiento de asistentes virtuales. Empresas multinacionales que deciden contrataciones a nivel global para todas sus filiales son claros ejemplos de clientes de este tipo de oferentes. La fortaleza está dada por la interconexión con otros servicios, la robustez, la diversidad de reportes y el soporte, pero como ocurre ante tanta estructura lo que se pierde es configurabilidad, flexibilidad y velocidad de adaptación.

\subsubsection{Definición del producto}

Considerando los puntos mencionados se deduce que la tecnología propuesta en la tesis debe enmarcarse en una solución llamada ``Software de hiperpersonalización para asistentes virtuales''.

Se instala localmente o en el entorno del cliente, actuando como un middleware entre las fuentes de datos y el motor del bot.

\textbf{Características principales:}
\begin{itemize}
    \item Observa el flujo de las conversaciones.
    \item Proporciona contexto adecuado basado en los datos disponibles.
    \item Asegura la hiperpersonalización sin comprometer la seguridad o privacidad de los datos.
\end{itemize}

\textbf{Ventajas:}
\begin{itemize}
    \item Reduce las preocupaciones de las empresas sobre compartir datos sensibles.
    \item Facilita la integración con sistemas existentes, como CRM o bases de datos internas.
    \item Puede ser adaptado a diferentes industrias y casos de uso.
\end{itemize}

Se venderá en licencias que se instalan en los entornos propios de los clientes. 